\section{Graph Theory}

\subsection{Floyd-Warshall}
All-pairs shortest paths in $O(V^3)$. Handles negative weights. Detects and propagates negative cycles: after main DP, if $m[k][k] < 0$ then all paths through $k$ become $-\infty$.
\begin{minted}{c++}
void floydWarshall(vector<vector<ll>>& m) {
  int n = m.size();
  const ll INF = 1LL << 62;
  for (int i = 0; i < n; i++)
    m[i][i] = min(m[i][i], 0LL);
  for (int k = 0; k < n; k++)
    for (int i = 0; i < n; i++)
      for (int j = 0; j < n; j++)
        if (m[i][k] != INF && m[k][j] != INF)
          m[i][j] = min(m[i][j],
            max(m[i][k]+m[k][j], -INF));
  for (int k = 0; k < n; k++)
    if (m[k][k] < 0)
      for (int i = 0; i < n; i++)
        for (int j = 0; j < n; j++)
          if (m[i][k]!=INF && m[k][j]!=INF)
            m[i][j] = -INF;
}
\end{minted}
\subsection{Bellman-Ford / SPFA}
Single-source shortest paths with negative edges. $O(VE)$. On $n$-th iteration, if distance still improves $\Rightarrow$ negative cycle. Sets distance to $-\infty$ for nodes reachable from negative cycle.
\begin{minted}{c++}
const ll LINF = LLONG_MAX / 2;
struct Edge { int a, b; ll w; };
struct BFNode { ll dist=LINF; int prev=-1; };
void bellmanFord(vector<BFNode>& nodes,
                 vector<Edge>& eds, int s) {
  nodes[s].dist = 0;
  int n = nodes.size();
  for (int i = 0; i < n; i++) {
    bool changed = false;
    for (Edge& e : eds) {
      if (nodes[e.a].dist == LINF) continue;
      ll d = nodes[e.a].dist + e.w;
      if (d < nodes[e.b].dist) {
        nodes[e.b].dist = (i < n-1 ? d : -LINF);
        nodes[e.b].prev = e.a;
        changed = true;
      }
    }
    if (!changed) break;
  }
  for (int i = 0; i < n; i++)
    for (Edge& e : eds)
      if (nodes[e.a].dist == -LINF)
        nodes[e.b].dist = -LINF;
}
\end{minted}
\subsection{Eulerian Path (Directed)}
Finds Eulerian path in directed graph. Returns empty vector if none exists. Requires in-degree == out-degree for all nodes (except start/end, diff $= 1$). \texttt{gr[u]} = list of \texttt{(v, edge\_id)}. \texttt{nedges} = total edges.
\begin{minted}{c++}
vector<int> eulerPath(
  vector<vector<pair<int,int>>>& gr,
  int nedges, int src=0) {
  int n = gr.size();
  vector<int> its(n), used(nedges), ret;
  vector<int> s = {src};
  while (!s.empty()) {
    int x = s.back(), &it = its[x];
    if (it == (int)gr[x].size()) {
      ret.push_back(x); s.pop_back();
    } else {
      auto [y,eid] = gr[x][it++];
      if (!used[eid]) { used[eid]=1; s.push_back(y); }
    }
  }
  if ((int)ret.size() != nedges+1) return {};
  reverse(ret.begin(), ret.end());
  return ret;
}
\end{minted}
\subsection{Eulerian Circuit}
Finds Eulerian circuit. Populates \texttt{path} with nodes in reverse order. Call \texttt{dfs(src)} and then \texttt{reverse(path.begin(), path.end())}. \texttt{visited} must be sized to total edges.
\begin{minted}{c++}
struct Edge { int to, edge_id; };
vector<vector<Edge>> g;
vector<bool> visited;
vector<int> path;
void dfs(int node) {
  while (!g[node].empty()) {
    auto e = g[node].back(); g[node].pop_back();
    if (visited[e.edge_id]) continue;
    visited[e.edge_id] = true;
    dfs(e.to);
  }
  path.push_back(node);
}
\end{minted}
\subsection{Graph Condense (Kosaraju SCC)}
Builds condensed DAG of strongly connected components. Two DFS passes: first on original graph, second on reversed graph. \texttt{comps} = number of SCCs. \texttt{comp[u]} = component id of node $u$. Key idea: DFS1 on original records finish-order; DFS2 on reverse assigns SCCs in reverse finish-order (Kosaraju's trick).
\begin{minted}{c++}
struct SCC {
  int n, comps;
  vector<vector<int>> adj, rev_adj, scc_adj;
  vector<int> order, comp;
  vector<bool> vis;
  SCC(int _n) : n(_n), comps(0), adj(_n), rev_adj(_n),
    comp(_n, -1), vis(_n, false) {}
  void add_edge(int u, int v) {
    adj[u].push_back(v); rev_adj[v].push_back(u);
  }
  void dfs1(int u) {
    vis[u] = true;
    for (int v : adj[u]) if (!vis[v]) dfs1(v);
    order.push_back(u);
  }
  void dfs2(int u, int id) {
    comp[u] = id;
    for (int v : rev_adj[u])
      if (comp[v] == -1) dfs2(v, id);
  }
  void build() {
    for (int i = 0; i < n; i++) if (!vis[i]) dfs1(i);
    reverse(order.begin(), order.end());
    for (int u : order)
      if (comp[u] == -1) dfs2(u, comps++);
    scc_adj.resize(comps);
    for (int u = 0; u < n; u++) {
      for (int v : adj[u]) {
        if (comp[u] != comp[v])
          scc_adj[comp[u]].push_back(comp[v]);
      }
    }
  }
};
\end{minted}
\subsection{2-SAT}
Each var $x_i$: node $2i$ (false), $2i+1$ (true). \texttt{either(a,b)} adds $(a \vee b)$. \texttt{setValue(x)} forces $x$ true. \texttt{solve()} assigns \texttt{values[i]=0/1}.
\begin{minted}{c++}
struct TwoSat {
  int N;
  vector<vector<int>> gr;
  vector<int> values;
  TwoSat(int n=0):N(n),gr(2*n){}
  int addVar(){gr.emplace_back();gr.emplace_back();
    return N++;}
  void either(int f,int j){
    f=max(2*f,-1-2*f);j=max(2*j,-1-2*j);
    gr[f].push_back(j^1);
    gr[j].push_back(f^1);
  }
  void setValue(int x){either(x,x);}
  vector<int> val,comp,z; int time=0;
  int dfs(int i){
    int low=val[i]=++time,x;
    z.push_back(i);
    for(int e:gr[i]) if(!comp[e])
      low=min(low,val[e]?:dfs(e));
    if(low==val[i]){
      do{x=z.back();z.pop_back();
        comp[x]=low;
        if(values[x>>1]==-1)values[x>>1]=x&1;
      }while(x!=i);
    }
    return val[i]=low;
  }
  bool solve(){
    values.assign(N,-1);
    val.assign(2*N,0);comp=val;
    for(int i=0;i<2*N;i++) if(!comp[i])dfs(i);
    for(int i=0;i<N;i++)
      if(comp[2*i]==comp[2*i+1]) return false;
    return true;
  }
};
// either(a,b)=a OR b; setValue(x)=x is true
// atMostOne({a,b,c})=at most one is true
\end{minted}
\subsection{Articulation Points \texorpdfstring{\&}{&} Bridges}
Finds cut vertices and cut edges in $O(N+M)$. Root is articulation if it has $\ge 2$ children in DFS tree. Edge $(u,v)$ is a bridge if $low[v] > tin[u]$. Key idea: $low[u]$ = earliest DFS entry time reachable from subtree of $u$ via at most one back-edge.
\begin{minted}{c++}
struct Articulation {
  int n, timer, ecnt;
  vector<vector<pair<int,int>>> g;
  vector<int> tin, low;
  vector<bool> is_art, is_bridge;
  vector<pair<int,int>> bridges;
  Articulation(int _n, int _m) : n(_n), timer(0),
    ecnt(0), g(_n), tin(_n, -1), low(_n, -1),
    is_art(_n, false), is_bridge(_m, false) {}
  void add_edge(int u, int v) {
    g[u].push_back({v, ecnt});
    g[v].push_back({u, ecnt++});
  }
  void dfs(int u, int peid = -1) {
    tin[u] = low[u] = timer++;
    int children = 0;
    for (auto [v, eid] : g[u]) {
      if (eid == peid) continue;
      if (tin[v] != -1) {
        low[u] = min(low[u], tin[v]);
      } else {
        children++;
        dfs(v, eid);
        low[u] = min(low[u], low[v]);
        if (low[v] >= tin[u] && peid != -1)
          is_art[u] = true;
        if (low[v] > tin[u]) {
          bridges.push_back({min(u,v), max(u,v)});
          is_bridge[eid] = true;
        }
      }
    }
    if (peid == -1 && children > 1) is_art[u] = true;
  }
  void build() {
    for (int i = 0; i < n; i++)
      if (tin[i] == -1) dfs(i);
  }
};
\end{minted}
\subsection{Boruvka's MST}
Time: $O(E \log V)$. Returns MST weight, or $-1$ if disconnected. Requires \texttt{dsu} template. Finds the minimum outgoing edge for each component simultaneously. Excellent for problems with implicit edges.
\begin{minted}{c++}
struct Edge { int u, v; ll w; };
ll boruvka(int n, const vector<Edge>& edges) {
  dsu d(n); ll ans = 0;
  while (d.cnt > 1) {
    vector<int> mn(n, -1);
    for (int i = 0; i < (int)edges.size(); ++i) {
      int u = d.find(edges[i].u);
      int v = d.find(edges[i].v);
      if (u == v) continue;
      if (mn[u] == -1 || edges[i].w < edges[mn[u]].w)
        mn[u] = i;
      if (mn[v] == -1 || edges[i].w < edges[mn[v]].w)
        mn[v] = i;
    }
    bool changed = false;
    for (int i = 0; i < n; ++i) {
      if (mn[i] == -1) continue;
      auto& e = edges[mn[i]];
      if (d.unite(e.u, e.v)) {
        ans += e.w; changed = true;
      }
    }
    if (!changed) break;
  }
  return d.cnt > 1 ? -1 : ans;
}
\end{minted}
\subsection{Biconnected Comp. \texorpdfstring{\&}{&} Bridge Tree}
Finds 2-edge-connected components (BCCs) -- maximal subgraphs with no bridges. Contract each BCC into a node; original bridges become edges of a tree (bridge tree).

How it works: First DFS finds bridges via $low[v] > tin[u]$. Second BFS floods each component avoiding bridges -- all nodes reachable without crossing a bridge get the same component id.
\begin{minted}{c++}
struct BCC {
  int n, timer, comps;
  vector<vector<int>> adj, tree;
  vector<int> tin, low, comp;
  vector<bool> is_bridge, vis;
  vector<pair<int, int>> edges;
  BCC(int _n) : n(_n), timer(0), comps(0), adj(_n),
    tin(_n, -1), low(_n, -1), comp(_n, -1),
    vis(_n, false) {}
  void add_edge(int u, int v) {
    adj[u].push_back(edges.size());
    adj[v].push_back(edges.size());
    edges.push_back({u, v});
    is_bridge.push_back(false);
  }
  void dfs(int u, int peid = -1) {
    vis[u] = true;
    tin[u] = low[u] = timer++;
    for (int id : adj[u]) {
      int v = edges[id].first ^
              edges[id].second ^ u;
      if (id == peid) continue;
      if (vis[v]) low[u] = min(low[u], tin[v]);
      else {
        dfs(v, id);
        low[u] = min(low[u], low[v]);
        if (low[v] > tin[u]) is_bridge[id] = true;
      }
    }
  }
  void build_tree() {
    for (int i = 0; i < n; i++) if (!vis[i]) dfs(i);
    vector<int> q;
    for (int i = 0; i < n; i++) {
      if (comp[i] != -1) continue;
      q.push_back(i); comp[i] = comps++;
      while (!q.empty()) {
        int u = q.back(); q.pop_back();
        for (int id : adj[u]) {
          int v = edges[id].first ^
                  edges[id].second ^ u;
          if (comp[v] == -1 && !is_bridge[id]) {
            comp[v] = comp[u]; q.push_back(v);
          }
        }
      }
    }
    tree.resize(comps);
    for (int i = 0; i < (int)edges.size(); i++) {
      if (is_bridge[i]) {
        int u = comp[edges[i].first],
            v = comp[edges[i].second];
        tree[u].push_back(v);
        tree[v].push_back(u);
      }
    }
  }
};
\end{minted}
\subsection{Binary Lifting}
Preprocess $O(N \log N)$, query $O(\log N)$. \texttt{lift(u,k)} returns $k$-th ancestor (or $-1$). \texttt{query(u,v)} returns LCA. Call \texttt{dfs(root,-1)} first. Key idea: $up[u][i] = 2^i$-th ancestor of $u$. Lift deeper node to same depth, then binary-search highest ancestor where paths diverge.
\begin{minted}{c++}
struct LCA {
  int n, LOG;
  vector<vector<int>> adj, up;
  vector<int> depth;
  LCA(int _n) : n(_n), LOG(__lg(_n) + 2), adj(_n),
    up(_n, vector<int>(LOG, -1)), depth(_n, 0) {}
  void add_edge(int u, int v) {
    adj[u].push_back(v); adj[v].push_back(u);
  }
  void dfs(int u, int p = -1, int d = 0) {
    up[u][0] = p;
    depth[u] = d;
    for (int i = 1; i < LOG; i++) {
      if (up[u][i - 1] != -1)
        up[u][i] = up[up[u][i - 1]][i - 1];
    }
    for (int v : adj[u])
      if (v != p) dfs(v, u, d + 1);
  }
  int lift(int u, int k) {
    for (int i = 0; i < LOG; i++) {
      if ((k >> i) & 1) u = up[u][i];
      if (u == -1) break;
    }
    return u;
  }
  int query(int u, int v) {
    if (depth[u] < depth[v]) swap(u, v);
    u = lift(u, depth[u] - depth[v]);
    if (u == v) return u;
    for (int i = LOG - 1; i >= 0; i--) {
      if (up[u][i] != up[v][i]) {
        u = up[u][i]; v = up[v][i];
      }
    }
    return up[u][0];
  }
};
\end{minted}
\subsection{Heavy Light Decomposition (HLD)}
Path queries on trees. Decomposes tree into $O(\log N)$ heavy paths. Call \texttt{dfs(root, 0)} then \texttt{decompose(root, root)}. Uses BIT internally -- swap with Lazy SegTree for range updates. \texttt{query(u,v)} returns sum on path $u \to v$. \texttt{addVal(idx, delta)} and \texttt{editVal(idx, value)} update node value.
\begin{minted}{c++}
struct HLD {
  int n, cur_pos = 1;
  vector<int> arr, parent, size, depth, heavy,
  head, pos;
  vector<vector<int>> adj;
  BIT<ll> segtree;
  HLD(int n) : n(n), arr(n + 1), parent(n + 1),
  size(n + 2), depth(n + 2), heavy(n + 2, 0),
  head(n + 2), pos(n + 2), adj(n + 2),
  segtree(BIT<ll>(n + 2)) {}
  void add_edge(int u, int v) {
    adj[u].push_back(v), adj[v].push_back(u);
  }
  void dfs(int u, int p) {
    parent[u] = p; size[u] = 1;
    int max_sz = 0;
    for (int v : adj[u]) {
      if (v == p) continue;
      depth[v] = depth[u] + 1;
      dfs(v, u);
      size[u] += size[v];
      if (size[v] > max_sz)
        max_sz = size[v], heavy[u] = v;
    }
  }
  void decompose(int u, int h) {
    head[u] = h; pos[u] = cur_pos++;
    segtree.add(pos[u], arr[u]);
    if (heavy[u]) decompose(heavy[u], h);
    for (int v : adj[u])
      if (v != parent[u] && v != heavy[u])
        decompose(v, v);
  }
  int query(int u, int v) {
    int res = 0;
    while (head[u] != head[v]) {
      if (depth[head[u]] < depth[head[v]])
        swap(u, v);
      res += segtree.sum(pos[head[u]], pos[u]);
      u = parent[head[u]];
    }
    if (depth[u] > depth[v]) swap(u, v);
    res += segtree.sum(pos[u], pos[v]);
    return res;
  }
  void addVal(int idx, int delta) {
    segtree.add(pos[idx], delta);
    arr[idx] += delta;
  }
  void editVal(int idx, int value) {
    segtree.add(pos[idx], value - arr[idx]);
    arr[idx] = value;
  }
};
\end{minted}
\subsection{Centroid Decomposition}
Recursively finds centroids to build $O(\log N)$ depth decomposition tree. Useful for path-counting queries (e.g., count paths of length $k$). Mark centroid removed, recurse on remaining subtrees.
\begin{minted}{c++}
vector<vector<int>> adj;
vector<bool> is_removed;
vector<int> subtree_size;
int get_subtree_size(int node, int parent = -1) {
  subtree_size[node] = 1;
  for (int child : adj[node]) {
    if (child == parent || is_removed[child]) continue;
    subtree_size[node] += get_subtree_size(child, node);
  }
  return subtree_size[node];
}
int get_centroid(int node, int tree_size,
                 int parent = -1) {
  for (int child : adj[node]) {
    if (child == parent || is_removed[child]) continue;
    if (subtree_size[child] * 2 > tree_size)
      return get_centroid(child, tree_size, node);
  }
  return node;
}
void build_centroid_decomp(int node = 0) {
  int centroid = get_centroid(node,
  get_subtree_size(node));
  // do something with centroid
  is_removed[centroid] = true;
  for (int child : adj[centroid]) {
    if (is_removed[child]) continue;
    build_centroid_decomp(child);
  }
}
\end{minted}
\subsection{DSU on Tree (Sack)}
$O(N \log N)$ offline subtree queries. Key idea: Process light children first (discard), then heavy child (keep). Re-add all light subtree elements at the end -- each element is re-added at most $O(\log N)$ times. The \texttt{keep} flag controls whether to clear after processing.
\begin{minted}{c++}
int timer = 0;
vector<int> sz, in, out, flat, big;
void dfs_sz(int u, int p) {
  sz[u] = 1, big[u] = 0, in[u] = ++timer,
  flat[timer] = u;
  for (int v : adj[u]) {
    if (v != p) {
      dfs_sz(v, u);
      sz[u] += sz[v];
      if (sz[v] > sz[big[u]]) big[u] = v;
    }
  }
  out[u] = timer;
}
void add(int u, int val) { /* update DS */ }
void dfs_sack(int u, int p, bool keep) {
  for (int v : adj[u])
    if (v != p && v != big[u]) dfs_sack(v, u, 0);
  if (big[u]) dfs_sack(big[u], u, 1);
  for (int v : adj[u])
    if (v != p && v != big[u])
      for (int i = in[v]; i <= out[v]; i++)
        add(flat[i], 1);
  add(u, 1);
  // answer queries for subtree u here
  if (!keep)
    for (int i = in[u]; i <= out[u]; i++)
      add(flat[i], -1);
}
// Init: sz/in/out/flat/big.assign(n+1, 0);
// dfs_sz(1, 0); dfs_sack(1, 0, 1);
\end{minted}
\subsection{Hopcroft-Karp}
Maximum bipartite matching. Left side size = $n$, right side size = $m$. $O(E \sqrt{V})$.

How it works: BFS builds layers of unmatched nodes, then DFS finds augmenting paths simultaneously (multi-augmenting). Each phase increases matching size by at least 1; at most $O(\sqrt{V})$ phases.
\begin{minted}{c++}
struct HK {
  int n, m;
  vector<vector<int>> adj;
  vector<int> L, R, d;
  HK(int n, int m) : n(n), m(m), adj(n), 
                     L(n, -1), R(m, -1), d(n) {}
  void add_edge(int u, int v) { adj[u].push_back(v); }
  bool bfs() {
    queue<int> q; bool ok = false;
    for (int i = 0; i < n; ++i) {
      if (L[i] == -1) d[i] = 0, q.push(i);
      else d[i] = -1;
    }
    while (!q.empty()) {
      int u = q.front(); q.pop();
      for (int v : adj[u]) {
        if (R[v] == -1) ok = true;
        else if (d[R[v]] == -1) 
          d[R[v]] = d[u] + 1, q.push(R[v]);
      }
    }
    return ok;
  }
  bool dfs(int u) {
    for (int v : adj[u]) {
      if (R[v] == -1 || 
      (d[R[v]] == d[u] + 1 && dfs(R[v]))) {
        L[u] = v; R[v] = u;
        return true;
      }
    }
    d[u] = -1; return false;
  }
  int max_matching() {
    int ans = 0;
    while (bfs()) for (int i = 0; i < n; ++i) 
      if (L[i] == -1 && dfs(i)) ans++;
    return ans;
  }
  vector<pair<int, int>> matching_edges() {
    vector<pair<int, int>> res;
    for (int i = 0; i < n; i++) 
      if (L[i] != -1) res.emplace_back(i, L[i]);
    return res;
  }
  pair<vector<int>, vector<int>> min_vertex_cover() {
    max_matching();
    vector<int> Z(n, 0), cvL, cvR;
    auto vis = [&](auto& self, int u) -> void {
      Z[u] = 1;
      for (int v : adj[u]) 
        if (R[v] != -1 && !Z[R[v]]) self(self, R[v]);
    };
    for (int i = 0; i < n; i++) 
      if (L[i] == -1) vis(vis, i);
    for (int i = 0; i < n; i++) 
      if (!Z[i]) cvL.push_back(i);
    for (int j = 0; j < m; j++) 
      if (R[j] != -1 && Z[R[j]]) cvR.push_back(j);
    return {cvL, cvR};
  }
};
\end{minted}
\subsection{Dinic's Algorithm (Max Flow)}
$O(V^2E)$ worst case, $O(E \sqrt{V})$ for unit capacities. Edge $(u,v)$ with capacity $c$ is added as TWO directed arcs. \texttt{e.flow} = current flow, \texttt{e.w} = capacity. Access residual capacity via \texttt{e.w - e.flow}. Use \texttt{directed=true} for one-way edges. After \texttt{max\_flow()}, \texttt{min\_cut(M)} returns the original edges crossing the min cut (both endpoints $\le M$; omit \texttt{M} for all) --- the last BFS leaves the source side in \texttt{lvl[v] >= 0} (sink side \texttt{-1}), so no extra pass is needed.

Key idea: BFS builds level graph; DFS sends blocking flow along it. "Dead-end" nodes are skipped via \texttt{ptr[u]} (current arc trick).
\begin{minted}{c++}
struct Dinic {
  struct edge { int to, rev; ll flow, w; };
  const ll inf = 1e18;
  int n, s, t;
  vector<vector<edge>> g;
  vector<int> lvl, ptr;
  vector<array<int,3>> edges;                 // {u, v, directed} of original edges
  Dinic(int n, int s, int t) : n(n), s(s), t(t), 
    g(n), lvl(n), ptr(n) {}
  void add_edge(int u, int v, ll w, bool directed = false) {
    if (u == v) return;
    edges.push_back({u, v, (int)directed});
    edge a = {v, (int)g[v].size(), 0, w};
    edge b = {u, (int)g[u].size(), 0, directed ? 0 : w};
    g[u].push_back(a), g[v].push_back(b);
  }
  bool bfs() {
    fill(lvl.begin(), lvl.end(), -1);
    lvl[s] = 0; queue<int> q; q.push(s);
    while (!q.empty()) {
      int u = q.front(); q.pop();
      for (auto &e : g[u]) {
        int v = e.to;
        if (e.flow < e.w && lvl[v] == -1) {
          lvl[v] = lvl[u] + 1; q.push(v);
        }
      }
    }
    return lvl[t] != -1;
  }
  ll dfs(int u, ll pushed) {
    if (pushed == 0 || u == t) return pushed;
    for (int &i = ptr[u]; i < (int)g[u].size(); ++i) {
      edge &e = g[u][i];
      if (e.w <= e.flow) continue;
      int v = e.to;
      if (lvl[u]+1 != lvl[v]) continue;
      ll tr = dfs(v, min(pushed, e.w - e.flow));
      if (tr == 0) continue;
      e.flow += tr; g[v][e.rev].flow -= tr;
      return tr;
    }
    return 0;
  }
  ll max_flow() {
    ll f = 0, pushed;
    while (bfs()) {
      fill(ptr.begin(), ptr.end(), 0);
      while ((pushed = dfs(s, inf)) > 0) f += pushed;
    }
    return f;                   // the last (failed) bfs left lvl[] = source side of the cut
  }
  // Min-cut edges among nodes with both endpoints <= M (call right after max_flow()).
  // lvl[v] >= 0  <=>  v on the source side;  lvl[v] == -1  <=>  sink side.
  vector<pair<int,int>> min_cut(int M = INT_MAX) {
    vector<pair<int,int>> ans;
    for (auto &[u, v, dir] : edges) if (max(u, v) <= M) {
      bool ru = lvl[u] >= 0, rv = lvl[v] >= 0;
      if (dir ? (ru && !rv) : (ru != rv)) ans.push_back({u, v});
    }
    sort(ans.begin(), ans.end());
    ans.erase(unique(ans.begin(), ans.end()), ans.end());
    return ans;
  }
};
\end{minted}
\subsection{Min Cost Max Flow}
Finds min-cost max-flow using potentials (Johnson's). $O(F \cdot E \log V)$. \texttt{addEdge(u,v,cap,cost)} adds edge $u \to v$ with capacity and cost. \texttt{min\_cost\_flow(s,t,K)} returns \texttt{\{actual\_flow, total\_cost\}} up to flow limit $K$. Returns \texttt{\{-1, -1\}} if a negative cycle is detected. Handles up to $\sim 10^3$ edges.
\begin{minted}{c++}
struct MCMF {
  struct edge { int v, rev; ll cap, cost; };
  int N;
  vector<vector<edge>> g;
  vector<ll> dist, pi;
  vector<int> edge_idx, par;
  const ll inf = 1e18;
  MCMF(int n):N(n),g(n),dist(n),pi(n),edge_idx(n),par(n){}
  void addEdge(int u,int v,ll cap,ll cost){
    if(u == v) return;
    g[u].push_back({v,(int)g[v].size(),cap,cost});
    g[v].push_back({u,(int)g[u].size()-1,0,-cost});
  }
  pair<ll,ll> min_cost_flow(int s,int t,ll K=inf){
    if(s == t) return {0,0};
    ll flow=0, cost=0;
    fill(pi.begin(),pi.end(),inf); pi[s]=0;
    for(int _=0;_<N;_++){
      bool any=false;
      for(int u=0;u<N;u++) if(pi[u]!=inf)
        for(auto& e:g[u])
          if(e.cap>0 && pi[u]+e.cost<pi[e.v])
            pi[e.v]=pi[u]+e.cost, any=true;
      if(!any) break;
    }
    for(int u=0;u<N;u++) if(pi[u]!=inf)
      for(auto& e:g[u])
        if(e.cap>0 && pi[u]+e.cost<pi[e.v]) return {-1,-1};
    while(flow < K){
      fill(dist.begin(),dist.end(),inf); dist[s]=0;
      priority_queue<pair<ll,int>, vector<pair<ll,int>>, greater<>> pq;
      pq.push({0,s});
      while(!pq.empty()){
        auto [d,u] = pq.top(); pq.pop();
        if(d>dist[u]) continue;
        for(int i=0;i<(int)g[u].size();i++){
          edge &e = g[u][i];
          if(e.cap>0){
            ll nd = dist[u]+e.cost+pi[u]-pi[e.v];
            if(nd<dist[e.v]){
              dist[e.v]=nd; par[e.v]=u; edge_idx[e.v]=i;
              pq.push({nd,e.v});
            }
          }
        }
      }
      if(dist[t]==inf) break;
      for(int i=0;i<N;i++) if(dist[i]<inf) pi[i]+=dist[i];
      ll f = K - flow;
      for(int u=t;u!=s;u=par[u]) f=min(f,g[par[u]][edge_idx[u]].cap);
      for(int u=t;u!=s;u=par[u]){
        edge &e = g[par[u]][edge_idx[u]];
        e.cap-=f; g[e.v][e.rev].cap+=f;
      }
      flow+=f; cost+=f*pi[t];
    }
    return {flow,cost};
  }
};
\end{minted}
\subsection{Link-Cut Tree}
Dynamic tree (link/cut). Supports adding/removing edges and connectivity queries, all in $O(\log N)$ amortized. 0-indexed nodes. \texttt{link(u,v)} -- add edge. \texttt{cut(u,v)} -- remove edge. \texttt{connected(u,v)} -- check if same tree. Each Node can store subtree data via \texttt{fix()}.
\begin{minted}{c++}
struct Node {
  Node *p = 0, *pp = 0, *c[2];
  bool flip = 0;
  Node() { c[0] = c[1] = 0; fix(); }
  void fix() {
    if (c[0]) c[0]->p = this;
    if (c[1]) c[1]->p = this;
  }
  void pushFlip() {
    if (!flip) return;
    flip = 0; swap(c[0], c[1]);
    if (c[0]) c[0]->flip ^= 1;
    if (c[1]) c[1]->flip ^= 1;
  }
  int up() { return p ? p->c[1] == this : -1; }
  void rot(int i, int b) {
    int h = i ^ b;
    Node *x = c[i], *y = b == 2 ? x : x->c[h];
    Node *z = b ? y : x;
    if ((y->p = p)) p->c[up()] = y;
    c[i] = z->c[i ^ 1];
    if (b < 2) {
      x->c[h] = y->c[h ^ 1];
      y->c[h ^ 1] = x;
    }
    z->c[i ^ 1] = this;
    fix(); x->fix(); y->fix();
    if (p) p->fix();
    swap(pp, y->pp);
  }
  void splay() {
    for (pushFlip(); p; ) {
      if (p->p) p->p->pushFlip();
      p->pushFlip(); pushFlip();
      int c1 = up(), c2 = p->up();
      if (c2 == -1) p->rot(c1, 2);
      else p->p->rot(c2, c1 != c2);
    }
  }
  Node* first() {
    pushFlip();
    return c[0] ? c[0]->first() : (splay(), this);
  }
};
struct LinkCut {
  vector<Node> node;
  LinkCut(int n) : node(n) {}
  void link(int u, int v) {
    makeRoot(&node[u]);
    node[u].pp = &node[v];
  }
  void cut(int u, int v) {
    Node *x = &node[u], *top = &node[v];
    makeRoot(top); x->splay();
    if (x->pp) x->pp = 0;
    else { x->c[0] = top->p = 0; x->fix(); }
  }
  bool connected(int u, int v) {
    Node* nu = access(&node[u])->first();
    return nu == access(&node[v])->first();
  }
  void makeRoot(Node* u) {
    access(u); u->splay();
    if (u->c[0]) {
      u->c[0]->p = 0;
      u->c[0]->flip ^= 1;
      u->c[0]->pp = u;
      u->c[0] = 0; u->fix();
    }
  }
  Node* access(Node* u) {
    u->splay();
    while (Node* pp = u->pp) {
      pp->splay(); u->pp = 0;
      if (pp->c[1])
        pp->c[1]->p = 0, pp->c[1]->pp = pp;
      pp->c[1] = u; pp->fix(); u = pp;
    }
    return u;
  }
};
\end{minted}
